\documentclass[11pt,letterpaper]{article}
\usepackage[T1]{fontenc}
\usepackage{lmodern}
\usepackage[left=.96in,right=.96in,top=.90in,bottom=.91in]{geometry}
\usepackage{amsmath,amssymb,amsthm,microtype}
\usepackage{xcolor}
\usepackage[colorlinks=true,linkcolor=blue!40!black,citecolor=blue!40!black,
  urlcolor=blue!40!black]{hyperref}
\hypersetup{pdftitle={A height criterion for irrational weight-three Lambert sums},
  pdfauthor={twinheterodox},pdfsubject={Number theory}}
\newtheorem{theorem}{Theorem}[section]
\newtheorem{corollary}[theorem]{Corollary}
\newtheorem{proposition}[theorem]{Proposition}
\newtheorem{lemma}[theorem]{Lemma}
\numberwithin{equation}{section}
\newcommand{\Q}{\mathbb Q}
\newcommand{\Z}{\mathbb Z}
\newcommand{\C}{\mathbb C}
\newcommand{\lp}{\log^{+}}
\title{\fontsize{16}{20}\selectfont A height criterion for irrational weight-three Lambert sums}
\author{twinheterodox}
\date{18 November 2026}
\begin{document}
\maketitle

\begin{abstract}
For an algebraic number \(b\) with \(|b|>1\), the weight-three Lambert
sum is \(Z(b)=\sum_{n\ge1}n^2/(b^n-1)\). We prove that \(Z(b)\)
lies outside \(\Q(b)\) whenever the Mahler measure \(M(b)\) satisfies
\(\log M(b)<(352\pi^2/2889)\,r\log|b|\), where \(r=1\) for a real
base and \(r=2\) otherwise. The criterion allows several conjugates
outside the unit disk and applies to nonintegral bases. In particular,
\(Z(3+\sqrt3)\notin\Q(\sqrt3)\), although both conjugates of the base
exceed one.
\end{abstract}

\section{Two expanding conjugates}

For \(|b|>1\), the weight-three Lambert sum
\[
                         Z(b)=\sum_{n=1}^{\infty}\frac{n^2}{b^n-1}
\]
converges absolutely. With \(q=1/b\), this is a \(q\)-analogue of
\(\zeta(3)\). It is irrational at every integer base \(b\ge2\)
\cite[Theorem~1.1]{Uplift}. For an algebraic base \(b\), every partial
sum lies in \(\Q(b)\). Can we show that the limit lies outside that field?

At \(b=2+\sqrt3\) the answer is yes. The other conjugate,
\(2-\sqrt3\), lies inside the unit disk, so the earlier algebraic-base
theorem gives \(Z(2+\sqrt3)\notin\Q(\sqrt3)\)
\cite[Theorem~1.2]{Previous}. Now increase the base by one. The field
stays \(\Q(\sqrt3)\) and the other conjugate crosses the unit circle to
\(3-\sqrt3>1\). We prove that \(Z(3+\sqrt3)\) still lies outside
\(\Q(\sqrt3)\) by quantifying how large the other conjugates can be.

For a specified algebraic base \(b\), write its primitive integer
minimal polynomial as \(a_0\prod_{j=1}^d(x-b_j)\), with \(a_0>0\).
Its \emph{Mahler measure} is
\[
                     M(b)=a_0\prod_{j=1}^d\max(1,|b_j|).
\]
The product counts each conjugate separately. In particular, a nonreal
base and its complex conjugate contribute two factors of \(|b|\).

\begin{theorem}\label{thm:height}
Let \(b\) be algebraic with \(|b|>1\). Put \(r=1\) for real \(b\) and
\(r=2\) for nonreal \(b\). If
\begin{equation}\label{eq:height}
              \frac{\log M(b)}{r\log|b|}<\frac{352\pi^2}{2889},
\end{equation}
then \(Z(b)\notin\Q(b)\).
\end{theorem}

Both roots \(3\pm\sqrt3\) have Mahler measure \(6\). To five decimal
places the ratio in \eqref{eq:height} is \(1.15273\) at \(3+\sqrt3\)
and \(7.54740\) at \(3-\sqrt3\), against a cutoff of \(1.20253\).
When the larger root is the base, the smaller root contributes
about \(15.27\%\) as much logarithmic growth as the larger one,
within the roughly \(20.25\%\) allowed by the criterion.

\begin{corollary}\label{cor:sqrt}
At \(b=3+\sqrt3\), the real number \(Z(b)\) lies outside \(\Q(\sqrt3)\).
\end{corollary}
\begin{proof}
The roots of \(x^2-6x+6\) are \(b\) and \(b'=3-\sqrt3\), so
\(M(b)=bb'=6\). Since \(\sqrt3>12/7\), we have \(b>33/7\) and
\(b'<9/7\). The inequality \(9^6<33\cdot7^5\) gives \((b')^6<b\).
Thus \(\log M(b)<(7/6)\log b\), and
\(7/6<352\pi^2/2889\) follows from \(\pi>31/10\).
\end{proof}

The same \(7/6\) bound gives a family of nonintegral examples. Let
\(b\) be the larger root of \(2x^2-mx+1\) for an integer \(m\ge129\).
Its discriminant \(m^2-8\) lies strictly between \((m-1)^2\) and
\(m^2\). The lower bound on \(m\) makes the polynomial negative at
\(2^6=64\), so \(2<b^{1/6}\). The other root is \(1/(2b)<1\),
and hence \(M(b)=2b<b^{7/6}\).
Theorem~\ref{thm:height} gives \(Z(b)\notin\Q(b)\) in this case as well.

In general, \(\log M(b)-r\log|b|\) adds together \(\log a_0\) and
the positive logarithms of the moduli of the remaining conjugates.
Both contributions are controlled by the degree in the base of certain
integer polynomials. That degree bounds their growth at the other
expanding conjugates and the power of \(a_0\) in the integer resultant
of Section~\ref{sec:places}. Section~\ref{sec:degree} proves that
each polynomial has the same degree in the base as its cyclotomic
clearing factor and identifies its leading coefficient.

Bundschuh and V\"a\"an\"anen used the height ratio in \eqref{eq:height}
for Lambert sums with periodic numerators \cite[\S2, (2)]{BV}.
For \(b=v/u\) with \(u<v\) coprime positive integers, \(M(b)=v\),
and \eqref{eq:height} becomes \(\log u/\log v<1-2889/(352\pi^2)\).
This is the weight-three rational-parameter range recorded in
\cite[\S7.1]{Uplift} and \cite[entry FUT-RATIONAL]{Companion},
where the degree bound is used for homogenization. The stated
restriction \(u/v\le1/2\) follows here from \(v>u^2\) and \(v\ge2\).

\section{The moments and their determinant}\label{sec:inputs}

We recall the polynomials whose values at conjugate bases will be multiplied.
Their two arguments have different roles. The first is the base \(t\).
The second is a variable \(X\) that can be set equal to the Lambert
sum at one embedding and left arbitrary at another.

First take \(t>1\) real. The Euler derivative \(\theta=x\,d/dx\)
sends \(x^j\) to \(jx^j\), so applying it twice supplies the weight
\(j^2\). The denominator comes from summing on a geometric sequence. Put
\[
 L_t(f)=\sum_{m\ge1}(\theta^2 f)(t^{-m}),\qquad
 L_t(x^j)=j^2\sum_{m\ge1}t^{-mj}=\frac{j^2}{t^j-1}\quad(j\ge1).
\]
Also \(L_t(1)=0\). For \(f\) holomorphic near \([0,t^{-1}]\) the sum
converges since \(\theta^2f(x)=xf'(x)+x^2f''(x)=O_f(x)\) near zero.

To see what happens at a pole, write
\(G(y)=\sum_{j\ge1}j^2y^j=y(1+y)/(1-y)^3\) for \(|y|<1\).
Summing the geometric expansions of the Lambert series gives
\(Z(t)=\sum_{m\ge1}G(t^{-m})\) by absolute convergence.
For the first pole \(f(x)=(1-x/t)^{-1}\), differentiating its power
series gives \(\theta^2 f(x)=G(x/t)\). The \(m\)th summand of \(L_t(f)\)
is therefore \(G(t^{-(m+1)})\), skipping the first term \(G(t^{-1})\).
The pole at \(t^a\) skips the first \(a\) terms in the same way, giving
\[
 L_t\big((1-x/t^a)^{-1}\big)=Z(t)-S_a(t),\qquad
 S_a(t)=\sum_{m=1}^aG(t^{-m}).
\]

Now treat \(t\) as an indeterminate. For \(K\ge1\) put \(r_a=t^a\)
for \(1\le a\le K\) and \(D(x)=\prod_{a=1}^K(1-x/r_a)\).
Division and partial fractions express \(p(x)/D(x)\) over \(\Q(t)\)
as a polynomial plus a linear combination of \((1-x/r_a)^{-1}\).
Define \(\mathcal L_{t,X}\) by giving the polynomial part the monomial
values above and each simple fraction the value \(X-S_a(t)\).
For real \(t>1\), setting \(X=Z(t)\) recovers \(L_t\).
These moments form the determinant
\[
 \Delta_{K,t}(X)=
 \det\left[\mathcal L_{t,X}\!\left(\frac{x^{i+j}}{D(x)}\right)\right]_{0\le i,j<K}.
\]
Its entries are affine in \(X\). The coefficients can have poles where
two nodes coincide or \(t^j=1\). Let \(C_K(t)\) be the monic product of
cyclotomic polynomials specified in \cite[(2.3)]{Previous}. It clears
the poles away from \(t=0\). Set
\[
 e_K=\frac{K^3-K}{2},\qquad d_K=\deg C_K,\qquad
 P_K(t,X)=C_K(t)t^{-e_K}\Delta_{K,t}(X).
\]
After division by \(t^{e_K}\), the result is still an integer polynomial.
The following four properties
are proved in \cite[Lemmas~2.1 and~4.1, Propositions~3.1 and~5.1]{Previous}.
We put \(\lp x=\log\max(1,x)\) for \(x\ge0\).

\begin{proposition}\label{prop:inputs}
Put \(\gamma=321/(64\pi^2)\) and \(\varepsilon=11/18-\gamma>0\).
The following hold.
\begin{enumerate}
\item \(P_K\in\Z[t,X]\), \(\deg_XP_K\le K\), and
      \(d_K\le\gamma K^3+O(K^2\log(2K))\).
\item For fixed \(|b|>1\),
      \(\log|P_K(b,Z(b))|\le-\varepsilon K^3\log|b|+
        O_b(K^2\log(2K))\), with \(\log 0=-\infty\).
\item For each fixed \(X_0\in\C\),
      \(\max_{|t|\le1}\lp|P_K(t,X_0)|=O_{X_0}(K^2\log^3(2K))\).
\item For each fixed \(|b|>1\) there are arbitrarily large \(K\)
      for which \(P_K(b,Z(b))\ne0\).
\end{enumerate}
\end{proposition}

The construction and the cyclotomic exponents come from the
integer-base weight-three determinant of uplift-universal-inc
\cite[\S\S1--5 and Appendix~A]{Uplift}. The cubic scale of decay for
Hankel determinants on geometric nodes already occurs in Zudilin's
work on the generalized \(q\)-logarithm \cite[\S4, (8)]{Zudilin}.
The estimates in Proposition~\ref{prop:inputs} use the complex-base
form of the weight-three determinant, including its bound on the
closed disk for every fixed second argument.

\section{An integer from all the conjugates}\label{sec:places}

Take \(b=3+\sqrt3\) and \(b'=3-\sqrt3\). If \(Z(b)=u+v\sqrt3\)
for rational \(u,v\), its field conjugate would be \(Y=u-v\sqrt3\).
The two factors in the norm of \(P_K(b,Z(b))\) would then be
\(P_K(b,Z(b))\) and \(P_K(b',Y)\). The series converges at both roots,
yet conjugation can change the behavior of a limit. For example, it
sends the sequence \((2-\sqrt3)^n\to0\) to \((2+\sqrt3)^n\to\infty\).
To cover the unknown \(Y\), we need a bound at \(b'\) for every fixed
second argument.
Proposition~\ref{prop:inputs}(2) gives decay for the first factor.
For the second, the degree in the \emph{base} controls growth from
the unit circle out to \(b'\).

The degree bound \(\deg_tP_K\le d_K\) from Lemma~\ref{lem:degree}
suffices. For fixed \(Y\), the reversed polynomial
\(u^{d_K}P_K(1/u,Y)\) is then
holomorphic at the origin. Apply the maximum modulus principle to it
at \(u=1/t\) for \(|t|>1\). Its modulus on the unit circle is bounded
by Proposition~\ref{prop:inputs}(3). Together with that same estimate
inside the disk this gives, for every \(t\in\C\),
\begin{equation}\label{eq:arch}
 \lp|P_K(t,Y)|\le d_K\lp |t|+O_Y(K^2\log^3(2K)).
\end{equation}
Thus the logarithmic cost at \(b'\) is at most \(d_K\log b'+o(K^3)\).
For a general base, the same degree controls the denominators.

\begin{proof}[Proof of Theorem~\ref{thm:height}, assuming Lemma~\ref{lem:degree}]
Suppose \(\xi=Z(b)\) belongs to \(\Q(b)\). With
\(d=[\Q(b):\Q]\), we may write \(\xi=g(b)/h\) for a fixed
\(g\in\Z[t]\) of degree at most \(d-1\) and a positive integer \(h\).
This representation specifies the conjugate value \(g(b_j)/h\) at
\emph{every} root \(b_j\) of the minimal polynomial \(f\).
Since \(\deg_X P_K\le K\),
\[
 Q_K(t)=h^K P_K\big(t,g(t)/h\big)\in\Z[t],\qquad
 \deg Q_K\le d_K+(d-1)K .
\]
Choose the unbounded sequence on which \(Q_K(b)=h^KP_K(b,\xi)\ne0\).
The irreducible \(f\) then has no common root with \(Q_K\).
Their resultant is a nonzero integer, as it is the determinant of
their integer Sylvester matrix. Its product expression is
\begin{equation}\label{eq:res}
 R_K=\operatorname{Res}(f,Q_K)
     =a_0^{\deg Q_K}h^{dK}\prod_{j=1}^d P_K\big(b_j,g(b_j)/h\big).
\end{equation}
Thus \(a_0\), which may exceed one for a nonintegral base, is raised
only to the degree of \(Q_K\). The fixed denominator \(h\) contributes
only \(O_{b,\xi}(K)\) to \(\log|R_K|\).

At \(b\) apply Proposition~\ref{prop:inputs}(2).
If \(b\) is nonreal, \(g(\bar b)/h=\bar\xi=Z(\bar b)\), giving a
second factor with the same decay. At each of the other \(d-r\)
roots use \eqref{eq:arch} with the fixed value \(Y=g(b_j)/h\).
In \eqref{eq:res} the leading coefficient \(a_0\) and these remaining
roots together contribute at most
\(d_K(\log M(b)-r\log|b|)+O_{b,\xi}(K^2\log^3(2K))\)
to the logarithm. Consequently
\begin{align*}
 0\le\log|R_K|
 &\le -r\varepsilon K^3\log|b|
       +d_K\big(\log M(b)-r\log|b|\big)+O_{b,\xi}(K^2\log^3(2K))\\
 &\le \left(\gamma\log M(b)-\tfrac{11}{18}r\log|b|\right)K^3+o(K^3).
\end{align*}
Here \(\log M(b)-r\log|b|\ge0\), so the upper bound for \(d_K\)
applies. Since \((11/18)/\gamma=352\pi^2/2889\), \eqref{eq:height}
makes the last coefficient strictly negative.
\end{proof}

\section{The degree in the base}\label{sec:degree}

\begin{lemma}\label{lem:degree}
For every \(K\ge1\), the degree of \(P_K(t,X)\) in \(t\) is \(d_K\),
and its coefficient of \(t^{d_K}\) is \((-1)^{K(K-1)/2}X^K\).
\end{lemma}

The remaining task is to see why the determinant contributes no
additional power of \(t\) after division by \(t^{e_K}\).
We let \(t\) tend to \(+\infty\) at fixed \(K\). The nodes \(r_a=t^a\)
separate, and the polynomials that interpolate at them make the
dependence on \(X\) diagonal.

\begin{proof}
Take real \(t\ge2\) and write \(L=L_t=\mathcal L_{t,Z(t)}\) as in
Section~\ref{sec:inputs}. Changing \(Z(t)\) to \(Y\) adds \(Y-Z(t)\) at each simple fraction
and leaves the polynomial moments fixed. The coefficient of the simple
fraction at \(r_a\) in \(p/D\) is \(w_a p(r_a)\), where
\[
           w_a=\prod_{c\ne a}(1-r_a/r_c)^{-1}.
\]
This follows by multiplying by \(1-x/r_a\) and evaluating at \(r_a\).
Thus the change in the functional on \(p/D\) is
\((Y-Z(t))\sum_a w_a p(r_a)\), even when division gives a polynomial
quotient.

Choose the degree-\((K-1)\) cardinal polynomials
\(\ell_a(x)=w_aD(x)/(1-x/r_a)\), characterized by
\(\ell_a(r_c)=\delta_{ac}\). In the matrix with entries
\(\mathcal L_{t,Y}(\ell_a\ell_c/D)\) the preceding change contributes
only to the diagonal. The full matrix is
\[
   (Y-Z(t))W+WHW,\qquad
   W=\operatorname{diag}(w_1,\ldots,w_K),\qquad
   H_{ac}=L\!\left(\frac{D(x)}{(1-x/r_a)(1-x/r_c)}\right).
\]
This also covers \(a=c\), when the fraction retains a simple pole.
The evaluation matrix from monomials to node values is
\(V=(r_a^i)_{1\le a\le K,\,0\le i<K}\). The coefficients of
the \(\ell_a\) are the columns of \(V^{-1}\), so returning to
monomials multiplies the determinant by \((\det V)^2\).
For each \(a<c\) the two factors \(1-r_a/r_c\) and \(1-r_c/r_a\)
have product \(-(r_c-r_a)^2/(r_a r_c)\). Hence
\[
   (\det V)^2\prod_a w_a=(-1)^{\binom K2}
        \left(\prod_a r_a\right)^{K-1}=(-1)^{\binom K2}t^{e_K}.
\]
Factoring \(W\) out of the matrix and multiplying by \(C_K(t)\) gives
\begin{equation}\label{eq:limit}
 P_K(t,Y)=(-1)^{\binom K2}C_K(t)\det(YI+E_t),
 \qquad E_t=HW-Z(t)I .
\end{equation}

For \(t\ge2\) the rational function defining \(H_{ac}\) has modulus
at most \(4\prod_{j\ge1}(1+2^{-j})\) on \(|x|\le1\).
Cauchy's estimates bound its first two derivatives on \(|x|\le1/2\).
Since \(\theta^2 f=xf'+x^2 f''\), its second Euler derivative is
\(O(|x|)\) there. Summing at \(t^{-m}\) yields \(H_{ac}=O(1/(t-1))\).
For fixed \(K\) the weights \(w_a\) stay bounded and \(Z(t)\to0\), so
\(E_t\to0\) entrywise. The determinant \(\det(YI+E_t)\) is monic
in \(Y\). Each lower coefficient is a sum of principal minors of
\(E_t\) of positive order and hence tends to zero.
By \eqref{eq:limit} the coefficient of \(Y^K\) in \(P_K(t,Y)\)
is exactly \((-1)^{\binom K2}C_K(t)\). Every lower coefficient is
\(o(t^{d_K})\) and is a polynomial in \(t\) by
Proposition~\ref{prop:inputs}(1), so it has degree less than \(d_K\)
whenever it is nonzero. Since \(C_K\) is monic of degree \(d_K\), the
coefficient of \(t^{d_K}\) in \(P_K(t,Y)\) is \((-1)^{\binom K2}Y^K\).
\end{proof}

\section*{Verification}
A collaborating AI model reviewed the mathematics and exposition of the
initial 18 November 2026 draft, taking Proposition~\ref{prop:inputs}
as a cited input. A separate AI model audited that draft for exposition.

\small
\begin{thebibliography}{9}
\setlength{\itemsep}{1pt}
\bibitem{BV} P.~Bundschuh and K.~V\"a\"an\"anen, Linear independence of certain
Lambert and allied series, \emph{Acta Arith.} \textbf{120} (2005), 197--209.
\href{https://doi.org/10.4064/aa120-2-7}{doi:10.4064/aa120-2-7}.
\bibitem{Previous} twinheterodox, The weight-three Lambert sum at the
golden ratio is irrational, \href{https://hexagonmath.org/2611.00623v1}{Hexagon 2611.00623v1} (2026).
\bibitem{Uplift} uplift-universal-inc, Irrationality of the Lambert-series
value \(\zeta_{1/p}(3)\), research draft (2026).
\href{https://github.com/uplift-universal-inc/qzeta_3-irrationality/blob/main/paper/main.pdf}{Manuscript}.
\bibitem{Companion} uplift-universal-inc, Results reserved for companion
articles, rational-parameter result and supporting degree bound (2026).
\href{https://github.com/uplift-universal-inc/qzeta_3-irrationality/blob/main/docs/FUTURE_RESULTS.md}{Companion inventory}.
\bibitem{Zudilin} W.~Zudilin, On the irrationality of generalized
\(q\)-logarithm, \emph{Res. Number Theory} \textbf{2} (2016), article~15.
\href{https://doi.org/10.1007/s40993-016-0042-x}{doi:10.1007/s40993-016-0042-x}.
\end{thebibliography}
\end{document}
